\section{Mathematical tools}
\label{sec:math-prelim}

\subsection{Wasserstein spaces}
\label{subsec:wasserstein}

Let $r\ge1$ be an integer, let $\mathscr P(\mathbb R^r)$ denote the Borel
probability measures on $\mathbb R^r$, and write $\nu_n\Rightarrow\nu$
for weak convergence (convergence against bounded continuous functions).

\begin{definition}[Wasserstein distance]
\label{definition:wasserstein-distance}
Write $\mathscr P_2(\mathbb R^r)$ for the measures $\mu$ satisfying
$\int\|x\|^2\,\mu(dx)<\infty$. For $\mu,\nu$ in this space, let
$\Pi(\mu,\nu)$ denote their couplings and define
\[
  W_2(\mu,\nu)
  :=\left(\inf_{\pi\in\Pi(\mu,\nu)}
     \int\|x-y\|^2\,\pi(dx,dy)\right)^{1/2}.
\]
Equivalently, the infimum is over joint laws of random vectors
$(X,Y)$ having marginals $\mu,\nu$, with objective
$(\mathbb E\|X-Y\|^2)^{1/2}$; see
\cite[Chapter~6]{villani2009optimal} and
\cite[Section~7.1]{ambrosio2008gradientflows}.
\end{definition}

\begin{definition}[Pseudo-Lipschitz functions]
\label{definition:pseudo-lipschitz}
A function $\psi:\mathbb R^r\to\mathbb R$ is
\emph{pseudo-Lipschitz of order $2$} if, for some $L<\infty$,
\[
  |\psi(x)-\psi(y)|
  \le L\|x-y\|(1+\|x\|+\|y\|),\qquad x,y\in\mathbb R^r.
\]
We denote this class by $\mathrm{PL}_2(\mathbb R^r;\mathbb R)$.
Its members are continuous and satisfy
$|\psi(x)|\le C(1+\|x\|^2)$ for some $C<\infty$.
\end{definition}

\begin{proposition}[Characterizations of $W_2$-convergence]
\label{proposition:wasserstein-convergence-characterization}
For deterministic $\nu_n,\bar\nu\in\mathscr P_2(\mathbb R^r)$,
the following are equivalent:
\begin{enumerate}
  \item $W_2(\nu_n,\bar\nu)\to0$.
  \item $\nu_n\Rightarrow\bar\nu$ and
  $\int\|x\|^2\,\nu_n(dx)\to\int\|x\|^2\,\bar\nu(dx)$.
  \item $\int\varphi\,d\nu_n\to\int\varphi\,d\bar\nu$ for every
  continuous $\varphi$ with $|\varphi(x)|\le C(1+\|x\|^2)$.
  \item $\int\psi\,d\nu_n\to\int\psi\,d\bar\nu$ for every
  $\psi\in\mathrm{PL}_2(\mathbb R^r;\mathbb R)$.
\end{enumerate}
If instead the $\nu_n$ are random and $\bar\nu$ is deterministic,
then $W_2(\nu_n,\bar\nu)\xrightarrow{\mathbb P}0$ if and only if
\[
  \int f\,d\nu_n\xrightarrow{\mathbb P}\int f\,d\bar\nu
  \quad
  \text{for every bounded Lipschitz }f, \text{ and}
  \quad
  \int\|x\|^2\,\nu_n(dx)
  \xrightarrow{\mathbb P}\int\|x\|^2\,\bar\nu(dx).
\]
\end{proposition}

\begin{proof}
The equivalence of the first three deterministic conditions follows from
\cite[Definition~6.8 and Theorem~6.9]{villani2009optimal},
\cite[Proposition~7.1.5]{ambrosio2008gradientflows}, and
\cite[Theorem~5.11]{santambrogio2015optimal}.
The third implies the fourth by the growth bound. Conversely,
bounded Lipschitz functions and $x\mapsto\|x\|^2$ belong to
$\mathrm{PL}_2$, so the fourth implies the second.

For random measures, necessity follows from the deterministic
characterization and the continuous mapping theorem
\cite[Lemma~5.3]{kallenberg2021foundations}.
For sufficiency, choose a countable convergence-determining family
of bounded Lipschitz functions; see
\cite[Proposition~2.2]{wormhille2009ergodic}.
From every subsequence, extract a further subsequence along which
all corresponding integrals and the second moments converge
almost surely, simultaneously
\cite[Lemmas~5.2 and~5.4]{kallenberg2021foundations}.
The deterministic characterization gives almost-sure $W_2$ convergence
along this further subsequence. The subsequence criterion
\cite[Lemma~5.2]{kallenberg2021foundations} completes the proof.
\end{proof}

\begin{lemma}[Lipschitz pushforwards]
\label{lemma:wasserstein-lipschitz-pushforward}
Let $T:E\to F$ be an $L$-Lipschitz map between Euclidean spaces.
For every $\mu,\nu\in\mathscr P_2(E)$,
\[
    W_2(T_\#\mu,T_\#\nu)\leq L W_2(\mu,\nu).
\]
In particular, coordinate projections cannot increase $W_2$.
\end{lemma}

\begin{proof}
For any coupling $\gamma$ of $\mu,\nu$, the measure
$(T,T)_\#\gamma$ couples $T_\#\mu,T_\#\nu$, and
\[
    \int\|T(x)-T(y)\|^2\,d\gamma
    \leq L^2\int\|x-y\|^2\,d\gamma.
\]
Taking infima proves the assertion.
\end{proof}

\subsection{Empirical measures and modes of convergence}
\label{subsec:empirical-measures}

\begin{definition}[Empirical measures]
\label{definition:empirical-measures}
For $A\in\mathbb R^{m\times r}$ with rows $A(i)$, set
\[
  \nu(A):=\frac1m\sum_{i=1}^m\delta_{A(i)}.
\]
For column vectors $a^1,\ldots,a^r\in\mathbb R^m$, write
$\nu(a^1,\ldots,a^r):=\nu([a^1\ \cdots\ a^r])$.
\end{definition}

\begin{corollary}[Convergence of empirical observables]
\label{corollary:wasserstein-empirical-observables}
If $A^{(m)}\in\mathbb R^{m\times r}$ are random and
$W_2(\nu(A^{(m)}),\bar\nu)\xrightarrow{\mathbb P}0$ for some
deterministic $\bar\nu\in\mathscr P_2(\mathbb R^r)$, then
\[
  \frac1m\sum_{i=1}^m\varphi(A^{(m)}(i))
  \xrightarrow{\mathbb P}\int\varphi\,d\bar\nu
\]
for every continuous $\varphi$ of at most quadratic growth.
The deterministic analogue holds with ordinary convergence.
\end{corollary}

\begin{proof}
Fix such a function $\varphi$.
Its integral is continuous for the $W_2$ topology by
Proposition~\ref{proposition:wasserstein-convergence-characterization}.
The deterministic assertion follows immediately, and the random
assertion follows by the continuous mapping theorem
\cite[Lemma~5.3]{kallenberg2021foundations}.
\end{proof}
